%==============================================================================
%  Detecção e Classificação de Lesões Mamárias em Ultrassonografia com YOLO11
%  Artigo científico — versão de submissão
%
%  Autocontido: usa apenas pacotes do TeX Live padrão.
%  Compile com pdfLaTeX (Overleaf: menu Recompile). Duas passadas resolvem
%  as referências cruzadas e a bibliografia.
%
%  >>> IMPORTANTE <<<
%  Todo número/afirmação em VERMELHO e entre colchetes — comando \pend — é um
%  PLACEHOLDER que depende dos resultados reais a serem rodados no Kaggle.
%  Não publicar enquanto houver \pend no texto. Ver checklist no README.
%==============================================================================
\documentclass[11pt,a4paper]{article}

\usepackage[T1]{fontenc}
\usepackage[utf8]{inputenc}
\usepackage[brazil]{babel}
\usepackage{lmodern}
\usepackage[margin=2.5cm]{geometry}
\usepackage{amsmath,amssymb}
\usepackage{graphicx}
\usepackage{booktabs}
\usepackage{tabularx}
\usepackage{xcolor}
\usepackage{microtype}
\usepackage[hidelinks]{hyperref}
\usepackage[round]{natbib}
\bibliographystyle{plainnat}

% --- Marcador de pendência (resultado real ainda não inserido) -----------------
\newcommand{\pend}[1][XX]{{\color{red}\textbf{[#1]}}}
% --- Nota de rodapé sem marcador (para a 1ª página) ---------------------------
\newcommand{\blfootnote}[1]{%
  \begingroup\renewcommand\thefootnote{}\footnote{#1}\addtocounter{footnote}{-1}\endgroup}

\title{\textbf{Detecção e classificação de lesões mamárias benignas e
malignas em ultrassonografia por meio de YOLO11: um estudo de prova de
conceito para auxílio ao diagnóstico}}

\author{%
Juan José Gouvêa Cardenas\thanks{Autor correspondente.} \and
Felipe Ramirez Pereira Botero%
\\[4pt]
\small \pend[Afiliação institucional --- preencher]\\
\small \texttt{\pend[e-mail de contato]}%
}
\date{\today}

\begin{document}
\maketitle

\blfootnote{\textbf{Aviso.} Este é um trabalho de pesquisa acadêmica em visão
computacional. O modelo descrito \textbf{não} é um dispositivo médico,
\textbf{não} realiza diagnóstico e \textbf{não} substitui a avaliação de um
profissional de saúde. Código e materiais:
\url{https://github.com/junowoz/deteccao-lesoes-mama-yolo}.}

%==============================================================================
\begin{abstract}
\noindent
O câncer de mama é a neoplasia mais incidente entre mulheres no Brasil, e o
diagnóstico tardio, agravado pela escassez de radiologistas na rede pública,
contribui para a mortalidade. A ultrassonografia é uma modalidade complementar
relevante, sobretudo em mamas densas, e sistemas de auxílio computacional (CAD)
podem apoiar a leitura dos exames. Este trabalho investiga o uso do detector de
objetos \emph{YOLO11s} para \emph{localizar} (caixa delimitadora) e
\emph{classificar} (benigna/maligna) lesões mamárias em imagens de ultrassom do
conjunto público BUSI. Diferentemente de boa parte da literatura sobre o BUSI,
que avalia a tarefa como classificação por imagem e reporta métricas
potencialmente infladas, adotamos um protocolo de avaliação adequado a
detectores: deduplicação por \emph{hash} antes da divisão dos dados (mitigando o
\emph{data leakage} documentado no BUSI), divisão estratificada
treino/validação/teste (70/15/15) com semente fixa e conjunto de teste cego, e
avaliação por \texttt{model.val()} com casamento por interseção sobre união
(IoU), reportando mAP@50, mAP@50--95 e precisão/revocação/F1 por classe, além da
taxa de falsos positivos em imagens normais. No conjunto de teste, o modelo
alcançou mAP@50 de \pend{} e mAP@50--95 de \pend{}, com revocação de \pend{} na
classe maligno. Os resultados são discutidos contra a faixa realista de
desempenho de \emph{detecção} em BUSI puro, e não contra índices de
\emph{classificação} de conjuntos aumentados, que não são comparáveis. O estudo
é uma prova de conceito reprodutível; suas limitações (defeitos conhecidos do
BUSI, origem monocêntrica, ausência de divisão por paciente e de validação
externa) são discutidas explicitamente.

\medskip
\noindent\textbf{Palavras-chave:} visão computacional; detecção de objetos; YOLO;
ultrassonografia mamária; aprendizado profundo; auxílio ao diagnóstico (CAD).
\end{abstract}

\renewcommand{\abstractname}{Abstract}
\begin{abstract}
\noindent
\textbf{(English)} Breast cancer is the most common cancer among women in
Brazil, and late diagnosis---worsened by the shortage of radiologists in the
public health system---contributes to mortality. Ultrasound is a relevant
complementary modality, especially for dense breasts, and computer-aided
detection/diagnosis (CAD) systems may support image reading. This work
investigates the \emph{YOLO11s} object detector to \emph{localise} (bounding
box) and \emph{classify} (benign/malignant) breast lesions in ultrasound images
from the public BUSI dataset. Unlike much of the BUSI literature, which treats
the task as image-level classification and reports potentially inflated metrics,
we adopt a detection-appropriate evaluation protocol: content-hash
deduplication before the split (mitigating BUSI's documented data leakage),
stratified train/validation/test split (70/15/15) with a fixed seed and a blind
test set, and \texttt{model.val()} evaluation with IoU matching, reporting
mAP@50, mAP@50--95 and per-class precision/recall/F1, plus the false-positive
rate on normal images. On the test set, the model reached mAP@50 of \pend{} and
mAP@50--95 of \pend{}. Results are a reproducible proof of concept; limitations
(known BUSI defects, single-centre origin, absence of patient-level split and
external validation) are explicitly discussed. This is not a medical device.

\medskip
\noindent\textbf{Keywords:} computer vision; object detection; YOLO; breast
ultrasound; deep learning; computer-aided diagnosis (CAD).
\end{abstract}

%==============================================================================
\section{Introdução}\label{sec:intro}

O câncer de mama é a neoplasia maligna mais incidente entre mulheres no Brasil.
O Instituto Nacional de Câncer (INCA) estimou cerca de 73.610 casos novos
por ano para o triênio 2023--2025, com taxa bruta de incidência de
66{,}54/100~mil mulheres \citep{inca2023}. A mortalidade permanece elevada e
está associada, entre outros fatores, ao diagnóstico em estágios avançados
\citep{inca2024}. No sistema público, o intervalo entre a suspeita e a
confirmação diagnóstica pode ultrapassar dois meses, com mediana de tempo até o
diagnóstico em torno de 70 dias e desempenho pior do que na rede privada
\citep{campos2022}. Esse cenário é agravado pela escassez de equipamentos e de
radiologistas, o que motiva o desenvolvimento de ferramentas de apoio à triagem
e à leitura de imagens.

A ultrassonografia mamária é uma modalidade complementar importante,
particularmente em mamas densas, nas quais a sensibilidade da mamografia é
reduzida. A padronização de laudos (por exemplo, pelo léxico BI-RADS) e a
variabilidade interobservador são desafios reconhecidos há décadas
\citep{pires2004}. Nesse contexto, sistemas de auxílio computacional ao
diagnóstico (\emph{Computer-Aided Detection/Diagnosis}, CAD) têm sido propostos
como ``segundo leitor'': não substituem o profissional, mas podem aumentar a
consistência e a eficiência da leitura. Evidência recente em rastreamento
mamográfico sugere que a leitura assistida por inteligência artificial pode ser
segura como apoio, sem aumento clinicamente relevante de achados de menor
interesse \citep{lang2023}.

Detectores de objetos em tempo real, como a família YOLO (\emph{You Only Look
Once}) \citep{redmon2016}, oferecem uma combinação atraente de velocidade e
desempenho de localização. Diferentemente de classificadores, que atribuem um
rótulo à imagem inteira, detectores \emph{localizam} a lesão (caixa
delimitadora) \emph{e} atribuem uma classe, o que se aproxima mais do fluxo de
trabalho da leitura assistida. Este trabalho investiga o uso do YOLO11
\citep{jocher2024,khanam2024} para detecção e classificação de lesões
benignas e malignas em ultrassom mamário, usando o conjunto público BUSI
\citep{aldhabyani2020}.

\paragraph{Enquadramento e escopo.} Trata-se de um estudo de \emph{prova de
conceito} em visão computacional. O sistema não diagnostica câncer: o
diagnóstico definitivo exige correlação clínica e confirmação histopatológica.
A saída do modelo deve ser entendida como \emph{achado sugestivo} de lesão de
aparência benigna ou maligna, em referência ao raciocínio BI-RADS, e nunca como
veredito. O trabalho não tem finalidade clínica, diagnóstica, terapêutica ou de
dispositivo médico.

\paragraph{Contribuições.} As contribuições deste trabalho são:
\begin{enumerate}
  \item Um \textbf{pipeline reprodutível} de detecção de lesões em ultrassom
  com YOLO11, derivando caixas delimitadoras a partir das máscaras de
  segmentação do BUSI, com semente fixa e divisão estratificada cega.
  \item Um \textbf{protocolo de avaliação adequado a detectores}
  (\texttt{model.val()} com IoU; mAP e métricas por classe; taxa de falsos
  positivos em imagens normais), em contraste com a prática comum de avaliar o
  BUSI como classificação por imagem.
  \item A \textbf{mitigação explícita do \emph{data leakage}} documentado no
  BUSI, por deduplicação de conteúdo antes da divisão, com discussão honesta das
  limitações do conjunto.
\end{enumerate}

%==============================================================================
\section{Trabalhos Relacionados}\label{sec:related}

\paragraph{Detecção e classificação em ultrassom mamário.} A detecção automática
de lesões em ultrassom mamário é um passo clássico de sistemas CAD.
\citet{yap2018} avaliaram abordagens de aprendizado profundo para a detecção de
lesões em ultrassom (incluindo uma rede totalmente convolucional por
transferência de aprendizado), estabelecendo um marco metodológico e apontando a
ausência de um conjunto comum como obstáculo à comparação entre métodos --- lacuna
que o BUSI \citep{aldhabyani2020} ajudaria a preencher. Trabalhos posteriores
com YOLO e variantes têm reportado desempenhos competitivos em ultrassom, embora
a comparabilidade entre estudos seja limitada por diferenças de conjunto, de
divisão e de protocolo de avaliação.

\paragraph{YOLO em imagem de mama.} Em mamografia, há evidência consistente de
viabilidade: \citet{prinzi2024} usaram YOLOv5s com transferência de aprendizado e
validação cruzada (mAP em torno de 0{,}835); \citet{anas2024} combinaram YOLOv5
com Mask R-CNN reportando acurácias elevadas em INbreast; e
\citet{mohammed2024} exploraram uma arquitetura YOLO multiescala com
pré-processamento por CLAHE. No Brasil, \citet{wurzel2022} avaliaram redes
convolucionais (VGG) para classificação em mamografia (CBIS-DDSM), com acurácia
em torno de 91\%. Esses trabalhos apoiam o uso de transferência de aprendizado e
de validação cruzada, mas em modalidade distinta (mamografia).

\paragraph{Cuidado com métricas não comparáveis.} Parte da literatura sobre o
BUSI reporta acurácias muito altas (94--99\%) ao tratar a tarefa como
\emph{classificação por imagem} e/ou ao treinar sobre dados fortemente
aumentados ou combinados. Por exemplo, \citet{jabeen2022} reportam acurácia
superior a 99\% sobre dados aumentados. Tais números medem uma tarefa diferente
(rótulo único por imagem, sem localização) e, frequentemente, sobre conjuntos
não diretamente comparáveis ao BUSI puro; não devem ser usados como linha de base
para \emph{detecção}. Esta distinção --- detecção (localização casada por IoU)
versus classificação por imagem --- orienta todo o protocolo de avaliação da
Seção~\ref{sec:metodos}.

\paragraph{Evolução das arquiteturas YOLO.} O YOLO11 \citep{jocher2024,khanam2024}
introduz blocos como o C3k2 e o C2PSA, com bom equilíbrio entre custo e
desempenho, e é a versão estável recomendada para uso prático. Versões
posteriores, como o YOLOv12 \citep{tian2025}, enfatizam mecanismos de atenção,
mas podem apresentar instabilidade de treino em conjuntos pequenos; o YOLO26
\citep{ultralytics2026}, com inferência sem supressão não máxima (NMS), é citado
como direção futura. Optou-se pelo YOLO11s por estabilidade e
reprodutibilidade em um conjunto de poucas centenas de imagens.

%==============================================================================
\section{Materiais e Métodos}\label{sec:metodos}

\subsection{Conjunto de dados BUSI}
Utilizou-se o \emph{Breast Ultrasound Images Dataset} (BUSI)
\citep{aldhabyani2020}, composto por 780 imagens de ultrassom adquiridas de 600
pacientes do sexo feminino (faixa etária de 25 a 75 anos), em formato PNG e
resolução média de aproximadamente $500\times500$ pixels, acompanhadas de
máscaras de segmentação. As imagens distribuídas dividem-se em 437 benignas, 210
malignas e 133 normais. As máscaras de segmentação (verdade de referência) são
usadas neste trabalho para derivar as caixas delimitadoras da tarefa de
detecção; imagens normais não possuem lesão e entram como exemplos negativos.

\subsection{Qualidade do conjunto e mitigação de \emph{data leakage}}
O BUSI apresenta defeitos documentados. \citet{pawlowska2023} relatam, entre
outros, cerca de 19\% de imagens duplicadas, além de imagens com marcadores de
medição, texto sobreposto, Color Doppler e estruturas não mamárias; estudos de
curadoria estimam que mais de 40\% das imagens apresentam algum defeito
significativo \citep{pawlowska2024}. Duplicatas distribuídas entre os conjuntos
de treino e teste causam \emph{data leakage} e inflam as métricas. Para mitigar
esse efeito, aplicou-se \textbf{deduplicação por hash MD5 do conteúdo} antes de
qualquer divisão, mantendo a primeira ocorrência de cada imagem. Após a
deduplicação, o conjunto resultante teve a seguinte composição:
\pend[treino/val/teste e contagens por classe --- preencher com a saída do notebook]{}.

\noindent\textit{Limitação estrutural.} O BUSI não fornece identificadores de
paciente; portanto, não é possível garantir divisão por paciente. Mesmo após a
remoção de duplicatas exatas, imagens de um mesmo paciente podem distribuir-se
entre conjuntos. Essa limitação é declarada na Seção~\ref{sec:limitacoes}.

\subsection{Geração de rótulos e divisão dos dados}
Para cada lesão (uma por máscara, contemplando os casos multilesão do BUSI),
a caixa delimitadora foi obtida pelo retângulo envolvente da região não nula da
máscara e normalizada para o formato YOLO ($x_c, y_c, w, h \in [0,1]$). As
classes de detecção são duas: \texttt{0 = benigno} e \texttt{1 = maligno}.
Imagens normais recebem arquivo de rótulo vazio (exemplo negativo/fundo).

A divisão foi \textbf{estratificada por classe} em treino (70\%), validação
(15\%) e teste (15\%), com semente fixa (\texttt{SEED = 42}) e conjunto de teste
mantido cego até a avaliação final. A divisão foi realizada \emph{antes} da
geração de rótulos e de qualquer aumento de dados.

\subsection{Modelo e treinamento}
Empregou-se o \textbf{YOLO11s} \citep{jocher2024} com pesos pré-treinados em COCO
(transferência de aprendizado), escolha adequada a um conjunto de poucas
centenas de imagens, reduzindo o risco de sobreajuste de variantes maiores. O
treinamento usou resolução de entrada $640\times640$, \emph{batch} de 16, até 100
épocas com \emph{early stopping} (\texttt{patience} = 25) e otimizador
selecionado automaticamente pelo Ultralytics.

\paragraph{Aumento de dados conservador.} Para preservar a anatomia do ultrassom,
restringiu-se o aumento a rotações leves (até $10^\circ$), espelhamento
horizontal ($p=0{,}5$), translação ($0{,}1$) e escala ($0{,}3$). Foram
\emph{desabilitados} cisalhamento, perspectiva e espelhamento vertical, por
introduzirem distorções geométricas pouco plausíveis na modalidade.

\subsection{Protocolo de avaliação}
A avaliação foi feita com \texttt{model.val()} (Ultralytics) sobre o conjunto de
teste, casando predições e verdade por interseção sobre união (IoU). Reportam-se:
\begin{itemize}
  \item \textbf{mAP@50} e \textbf{mAP@50--95} (médias sobre as classes);
  \item \textbf{precisão, revocação e F1 por classe}, sendo a revocação
  equivalente à sensibilidade --- métrica prioritária na classe maligno;
  \item \textbf{taxa de falsos positivos em imagens normais}: proporção de
  imagens sem lesão que receberam ao menos uma detecção espúria.
\end{itemize}
Adicionalmente, registram-se a matriz de confusão e a curva
\emph{precision--recall} geradas automaticamente pelo Ultralytics.

\subsection{Reprodutibilidade}
Toda operação aleatória usa semente fixa. O ambiente de execução
(Kaggle, GPU NVIDIA Tesla T4) e as versões dos pacotes são registrados
programaticamente pelo notebook
(\texttt{ultralytics}~\pend[versão]{}, \texttt{torch}~\pend[versão]{},
Python~\pend[versão]{}). O código, o notebook e as instruções de reprodução
estão disponíveis no repositório público (Seção~\ref{sec:disponibilidade}).

%==============================================================================
\section{Resultados}\label{sec:resultados}

\noindent\textbf{\textcolor{red}{Nota aos autores:} todos os valores em vermelho
abaixo são placeholders. Substitua-os pelos números do arquivo
\texttt{resultados\_experimentais.json} / \texttt{tabela\_resultados.md} gerado
pelo notebook após o re-treino. Remova esta nota antes de submeter.}

\subsection{Composição dos conjuntos após deduplicação}
A Tabela~\ref{tab:splits} resume o número de imagens por conjunto e classe após a
deduplicação e a divisão estratificada.

\begin{table}[h]
\centering
\caption{Composição dos conjuntos após deduplicação e divisão estratificada
(70/15/15). Valores a preencher com a saída do notebook.}
\label{tab:splits}
\begin{tabularx}{0.85\textwidth}{Xcccc}
\toprule
Conjunto & Benigno & Maligno & Normal & Total \\
\midrule
Treino     & \pend{} & \pend{} & \pend{} & \pend{} \\
Validação  & \pend{} & \pend{} & \pend{} & \pend{} \\
Teste      & \pend{} & \pend{} & \pend{} & \pend{} \\
\midrule
Total      & \pend{} & \pend{} & \pend{} & \pend{} \\
\bottomrule
\end{tabularx}
\end{table}

\subsection{Métricas de detecção no conjunto de teste}
A Tabela~\ref{tab:global} apresenta as métricas globais e a
Tabela~\ref{tab:classe} as métricas por classe, no conjunto de teste cego.

\begin{table}[h]
\centering
\caption{Métricas globais de detecção no conjunto de teste. Valores a preencher.}
\label{tab:global}
\begin{tabularx}{0.8\textwidth}{Xcccc}
\toprule
& mAP@50 & mAP@50--95 & Precisão & Revocação \\
\midrule
YOLO11s & \pend{} & \pend{} & \pend{} & \pend{} \\
\bottomrule
\end{tabularx}
\end{table}

\begin{table}[h]
\centering
\caption{Métricas por classe no conjunto de teste. Valores a preencher.}
\label{tab:classe}
\begin{tabularx}{0.92\textwidth}{Xccccc}
\toprule
Classe & Precisão & Revocação & F1 & AP@50 & AP@50--95 \\
\midrule
Benigno & \pend{} & \pend{} & \pend{} & \pend{} & \pend{} \\
Maligno & \pend{} & \pend{} & \pend{} & \pend{} & \pend{} \\
\bottomrule
\end{tabularx}
\end{table}

\subsection{Falsos positivos em imagens normais}
Das \pend[N]{} imagens normais do conjunto de teste, \pend[k]{}
($\pend[p]{}\%$) geraram ao menos uma detecção espúria. Esse indicador estima a
propensão do modelo a falsos positivos em casos sem lesão, relevante para um
cenário de triagem.

\subsection{Análise qualitativa}
A Figura~\ref{fig:exemplos} ilustra detecções representativas no conjunto de
teste (casos benignos, malignos e normais). \pend[Inserir figura
\texttt{inference\_results/} e a matriz de confusão \texttt{confusion\_matrix.png}
exportadas pelo notebook.]{}

\begin{figure}[h]
\centering
\fbox{\parbox[c][4cm][c]{0.8\textwidth}{\centering \pend[Espaço reservado para a
figura de exemplos de detecção e/ou matriz de confusão.]{}}}
\caption{Exemplos qualitativos de detecção no conjunto de teste. A inserir após o
re-treino.}
\label{fig:exemplos}
\end{figure}

%==============================================================================
\section{Discussão}\label{sec:discussao}

Os resultados devem ser interpretados como uma \emph{prova de conceito} de
detecção em ultrassom mamário, e não como evidência de desempenho clínico.
\pend[Discutir os valores reais obtidos:]{} em particular, comparar o mAP@50 e o
mAP@50--95 com a faixa realista de \emph{detecção} em BUSI puro relatada na
literatura, e \textbf{evitar} a comparação direta com índices de
\emph{classificação} de 94--99\% obtidos sobre dados aumentados/combinados
\citep{jabeen2022}, que medem tarefa distinta e não são comparáveis
(Seção~\ref{sec:related}).

A revocação na classe maligno é a métrica de maior interesse em um cenário de
apoio à triagem, pois falsos negativos têm custo clínico mais elevado.
\pend[Comentar o valor de revocação obtido para maligno e o equilíbrio
precisão/revocação.]{} A taxa de falsos positivos em imagens normais
complementa essa leitura, indicando quantas vezes o modelo ``alarmaria''
indevidamente em casos sem lesão.

É importante notar que parte expressiva do desempenho --- e de sua eventual
fragilidade --- decorre da \textbf{qualidade do conjunto}. A deduplicação prévia
tende a \emph{reduzir} métricas em relação a pipelines que treinam e avaliam
sobre o BUSI bruto, justamente por eliminar o vazamento que as inflava. Assim,
números mais modestos, porém honestos, são preferíveis a índices elevados
contaminados por \emph{leakage}.

%==============================================================================
\section{Limitações}\label{sec:limitacoes}

\begin{itemize}
  \item \textbf{Defeitos do conjunto.} Mesmo após a deduplicação exata, o BUSI
  contém imagens com marcadores, texto, Doppler e estruturas não mamárias
  \citep{pawlowska2023,pawlowska2024}, que não foram integralmente removidas
  neste estudo.
  \item \textbf{Ausência de divisão por paciente.} O BUSI não fornece IDs de
  paciente; pode haver imagens de um mesmo paciente em conjuntos diferentes.
  \item \textbf{Origem monocêntrica e tamanho reduzido.} O conjunto provém de
  uma única fonte e tem poucas centenas de imagens, o que limita a
  generalização. Uma única partição treino/teste tende a apresentar variância
  alta; validação cruzada é desejável (Seção~\ref{sec:conclusao}).
  \item \textbf{Sem validação externa.} Não foi avaliada a transferência para
  outros equipamentos/populações (por exemplo, conjuntos curados externos
  \citep{pawlowska2024}).
  \item \textbf{Escopo não clínico.} O trabalho não constitui validação clínica
  nem evidência de utilidade diagnóstica; não é um dispositivo médico.
\end{itemize}

%==============================================================================
\section{Conclusão e Trabalhos Futuros}\label{sec:conclusao}

Apresentou-se uma prova de conceito reprodutível de detecção e classificação de
lesões mamárias benignas e malignas em ultrassom com YOLO11s, com ênfase em rigor
metodológico: deduplicação prévia, divisão estratificada cega e avaliação
adequada a detectores. \pend[Resumir a conclusão à luz dos resultados reais
(por exemplo: ``o modelo alcançou mAP@50 de XX, demonstrando viabilidade da
abordagem como ferramenta de apoio, ressalvadas as limitações de dados'').]{}

Como trabalhos futuros, destacam-se: (i) validação cruzada $k$-fold para reduzir
a variância das estimativas; (ii) uso de versões curadas do BUSI e validação
externa em conjuntos independentes; (iii) tratamento explícito do
desbalanceamento entre classes para melhorar a revocação em maligno; (iv)
comparação entre variantes do YOLO (n/s/m) e contra uma linha de base de
classificação; e (v) avaliação de versões mais recentes com inferência sem NMS
\citep{ultralytics2026}.

%==============================================================================
\section*{Declarações}

\subsection*{Considerações éticas e de uso}
Trabalho de pesquisa em visão computacional sobre dados públicos e anonimizados
(BUSI). Não envolve coleta de dados de pacientes pelos autores nem intervenção
clínica. O sistema descrito não é um dispositivo médico, não realiza diagnóstico
e não substitui avaliação profissional; eventual uso assistivo exigiria
validação clínica e aprovação regulatória apropriadas.

\subsection*{Disponibilidade de dados e código}\label{sec:disponibilidade}
O conjunto BUSI é público \citep{aldhabyani2020}. O código, o notebook de
treinamento/avaliação e as instruções de reprodução estão disponíveis em
\url{https://github.com/junowoz/deteccao-lesoes-mama-yolo}. Os artefatos de
resultados (\texttt{resultados\_experimentais.json}, figuras e tabelas) são
gerados pelo notebook.

\subsection*{Contribuições dos autores (CRediT)}
\pend[Revisar/ajustar conforme a realidade do trabalho.]{}
\textbf{J. J. Gouvêa Cardenas:} concepção, metodologia, software,
experimentação, redação. \textbf{F. Ramirez Pereira Botero:} metodologia,
validação, análise, revisão e edição.

\subsection*{Conflitos de interesse}
Os autores declaram não haver conflitos de interesse.

\subsection*{Financiamento}
\pend[Declarar financiamento ou ``Os autores não receberam financiamento
específico para este trabalho''.]{}

%==============================================================================
\begin{thebibliography}{99}
\setlength{\itemsep}{2pt}

\bibitem[Al-Dhabyani et al.(2020)]{aldhabyani2020}
AL-DHABYANI, W.; GOMAA, M.; KHALED, H.; FAHMY, A.
Dataset of breast ultrasound images.
\emph{Data in Brief}, v. 28, art. 104863, 2020.
DOI: 10.1016/j.dib.2019.104863.

\bibitem[Anas et al.(2024)]{anas2024}
ANAS, M.; HAQ, I. U.; HUSNAIN, G.; JAFFERY, S. A. F.
Advancing breast cancer detection: enhancing YOLOv5 network for accurate
classification in mammogram images.
\emph{IEEE Access}, v. 12, p. 16474--16488, 2024.
DOI: 10.1109/ACCESS.2024.3358686.

\bibitem[Campos et al.(2022)]{campos2022}
CAMPOS, A. A. L. et al.
Time to diagnosis and treatment for breast cancer in public and private health
services.
\emph{Revista Gaúcha de Enfermagem}, v. 43, e20210103, 2022.
DOI: 10.1590/1983-1447.2022.20210103.en.

\bibitem[INCA(2022)]{inca2023}
INSTITUTO NACIONAL DE CÂNCER (INCA).
\emph{Estimativa 2023: incidência de câncer no Brasil}.
Rio de Janeiro: INCA, 2022.

\bibitem[INCA(2024)]{inca2024}
INSTITUTO NACIONAL DE CÂNCER (INCA).
\emph{Controle do câncer de mama no Brasil: dados e números 2024}.
Rio de Janeiro: INCA, 2024.

\bibitem[Jabeen et al.(2022)]{jabeen2022}
JABEEN, K. et al.
Breast cancer classification from ultrasound images using probability-based
optimal deep learning feature fusion.
\emph{Sensors}, v. 22, n. 3, art. 807, 2022.
DOI: 10.3390/s22030807.

\bibitem[Jocher e Qiu(2024)]{jocher2024}
JOCHER, G.; QIU, J.
\emph{Ultralytics YOLO11}. Software, v. 11.0.0, 2024.
Disponível em: \url{https://github.com/ultralytics/ultralytics}.

\bibitem[Khanam e Hussain(2024)]{khanam2024}
KHANAM, R.; HUSSAIN, M.
YOLOv11: an overview of the key architectural enhancements.
\emph{arXiv}:2410.17725, 2024.
DOI: 10.48550/arXiv.2410.17725.

\bibitem[Läng et al.(2023)]{lang2023}
LÄNG, K. et al.
Artificial intelligence-supported screen reading versus standard double reading
in the Mammography Screening with Artificial Intelligence trial (MASAI): interim
safety analysis.
\emph{The Lancet Oncology}, v. 24, n. 8, p. 936--944, 2023.
DOI: 10.1016/S1470-2045(23)00298-X.

\bibitem[Mohammed e Ekmekci(2024)]{mohammed2024}
MOHAMMED, A. D.; EKMEKCI, D.
Breast cancer diagnosis using YOLO-based multiscale parallel CNN and flattened
threshold Swish.
\emph{Applied Sciences}, v. 14, n. 7, art. 2680, 2024.
DOI: 10.3390/app14072680.

\bibitem[Pawlowska et al.(2023)]{pawlowska2023}
PAWLOWSKA, A.; KARWAT, P.; ZOLEK, N.
Letter to the editor. Re: ``Dataset of breast ultrasound images''.
\emph{Data in Brief}, v. 48, art. 109247, 2023.
DOI: 10.1016/j.dib.2023.109247.

\bibitem[Pawlowska et al.(2024)]{pawlowska2024}
PAWLOWSKA, A. et al.
A curated benchmark dataset for ultrasound based breast lesion analysis (BrEaST).
\emph{Scientific Data}, v. 11, 2024.
DOI: 10.1038/s41597-024-02984-z.

\bibitem[Pires et al.(2004)]{pires2004}
PIRES, S. R.; MEDEIROS, R. B.; SCHIABEL, H.
Banco de imagens mamográficas para treinamento na interpretação de imagens
digitais.
\emph{Radiologia Brasileira}, v. 37, n. 4, p. 239--244, 2004.
DOI: 10.1590/S0100-39842004000400006.

\bibitem[Prinzi et al.(2024)]{prinzi2024}
PRINZI, F.; INSALACO, M.; ORLANDO, A.; GAGLIO, S.; VITABILE, S.
A YOLO-based model for breast cancer detection in mammograms.
\emph{Cognitive Computation}, v. 16, p. 107--120, 2024.
DOI: 10.1007/s12559-023-10189-6.

\bibitem[Redmon et al.(2016)]{redmon2016}
REDMON, J.; DIVVALA, S.; GIRSHICK, R.; FARHADI, A.
You Only Look Once: unified, real-time object detection.
In: \emph{IEEE Conference on Computer Vision and Pattern Recognition (CVPR)},
2016. p. 779--788.
DOI: 10.1109/CVPR.2016.91.

\bibitem[Tian et al.(2025)]{tian2025}
TIAN, Y.; YE, Q.; DOERMANN, D.
YOLOv12: attention-centric real-time object detectors.
\emph{arXiv}:2502.12524, 2025.
DOI: 10.48550/arXiv.2502.12524.

\bibitem[Ultralytics(2026)]{ultralytics2026}
ULTRALYTICS.
\emph{Ultralytics YOLO26 documentation}. 2025--2026.
Disponível em: \url{https://docs.ultralytics.com/models/yolo26}.

\bibitem[Wurzel e Martins(2022)]{wurzel2022}
WURZEL, P.; MARTINS, M. O.
Classificação de imagens de mamografia com Machine Learning no auxílio de
diagnósticos de câncer de mama.
\emph{Disciplinarum Scientia: Naturais e Tecnológicas}, v. 23, n. 2, p. 1--17,
2022.
DOI: 10.37779/nt.v23i2.4140.

\bibitem[Yap et al.(2018)]{yap2018}
YAP, M. H. et al.
Automated breast ultrasound lesions detection using convolutional neural
networks.
\emph{IEEE Journal of Biomedical and Health Informatics}, v. 22, n. 4,
p. 1218--1226, 2018.
DOI: 10.1109/JBHI.2017.2731873.

\end{thebibliography}

\end{document}
